\documentclass[11pt]{article}
\usepackage[margin=1in]{geometry}
\usepackage{booktabs}
\usepackage{graphicx}
\usepackage{amsmath}
\usepackage{multirow}
\usepackage{url}
\usepackage{hyperref}
\usepackage{xcolor}
\newcommand{\todo}[1]{\textcolor{red}{[TODO: #1]}}
% Every number below must match docs/paper/EVIDENCE_MAP.md. Do not edit a
% number here without updating the map first.

\title{The Promotion Gap: Verified End-to-End Deployment of\\
Agent-Generated GPU Optimizations for Video Diffusion Inference}
\author{Aryan \todo{full author list and affiliations}\\
\todo{acknowledge agent tooling in Acknowledgments, not authorship}}
\date{August 2026}

\begin{document}
\maketitle

\begin{abstract}
Coding agents now routinely produce GPU kernels that pass correctness
suites and report large isolated speedups. We show that, for production
video diffusion inference, these two signals are nearly uncorrelated with
shipped value, and we present a promotion pipeline that measures
optimizations the way users experience them: full-generation latency and
output quality on declared workloads, under fail-closed verification. Run
against FastVideo on NVIDIA GB200, the pipeline surfaced a systematic
\emph{promotion gap}: three fused Wan~2.1 operators with $8.6$--$10.6\times$
isolated speedups left 50-step generation latency unchanged; four LTX
candidates covering $6.2\%$ of end-to-end time could never exceed
$1.0064\times$; and the delivery mechanism itself charged $3.104$\,ms per
call to realize a $124$\,\textmu s saving---a $25{:}1$ dispatch tax
consuming $32.4\%$ of end-to-end time before we eliminated it. Independent
validation caught agent-generated kernels that were bit-exact yet compiled
one call's memory layout in as constants, and dispatch code that silently
bypassed compiler guards to report an unfair $2.557\times$. Extending the
gates with tiered fidelity contracts (bitwise, numeric, perceptual), we
evaluate approximate optimizations honestly: two published attention
methods lose to a tuned FlashAttention baseline in this regime
($0.80\times$ and $0.52\times$) while structural and perceptual quality
metrics systematically disagree about them. One transformer artifact was
promoted end to end ($1.086\times$ median, byte-identical frames, $6{,}143$
dispatches, zero fallbacks). Every number in this paper traces to a
committed, machine-generated measurement record. We argue this promotion
layer---not kernel generation itself---is the missing infrastructure for
the era of agent-generated performance code.
\end{abstract}

\section{Introduction}

The cost of producing plausible GPU kernels has collapsed. Coding agents
iterate on Triton and CUDA faster than humans can review the output, and
benchmark suites report that they increasingly succeed: generated kernels
pass correctness tests and beat reference implementations by comfortable
margins. Evaluation practice, however, has not kept up with generation
practice. The dominant benchmarks score a kernel in isolation---correct
outputs on sampled inputs, speedup against a reference on fixed
shapes~\cite{kernelbench2025,kernelgenbench2026}. These are the two
signals an optimizing agent is trained and prompted to maximize, and we
show that in a production system they are nearly uncorrelated with
delivered value.

We report on MotionKernel, a system that holds agent-generated
optimizations to a different standard: an optimization counts only if it
improves \emph{full-generation} latency on a \emph{declared workload}
without violating a \emph{declared fidelity contract}, measured end to
end in the target framework with the artifact actually dispatching. The
pipeline is fail-closed at every stage---capture, ranking, spec
generation, search, validation, packaging, dispatch, and promotion each
reject by default and must be affirmatively passed. We applied it to
video diffusion inference (FastVideo~\cite{fastvideo} running Wan
2.1~\cite{wan2025} and LTX-2 on NVIDIA GB200), a domain whose cost
structure---hundreds of transformer invocations per generation, quadratic
attention over $10^4$--$10^5$ tokens, a heavy VAE---makes the distance
between operator microbenchmarks and user-visible latency unusually
stark.

That distance is the paper's central finding, which we call the
\emph{promotion gap}. Three fused Wan operators achieved
$8.6$--$10.6\times$ isolated speedups across their production shape
corpora and left 50-step generation latency unchanged
(Table~\ref{tab:wan}). Four validated LTX candidates with isolated
speedups of $1.11$--$1.12\times$ covered $6.20\%$ of end-to-end time in
aggregate: even at zero overhead and perfect fidelity, no combination
could exceed $1.0064\times$ (Table~\ref{tab:ltx-arith}). The delivery
mechanism itself initially charged $3.104$\,ms per call to realize a
$124$\,\textmu s saving---a $25{:}1$ dispatch tax that consumed $32.4\%$
of end-to-end time at 384 calls per generation
(Section~\ref{sec:dispatch}). And when we extended the gates to admit
approximate methods under explicit perceptual budgets, two published
attention accelerators lost to a tuned FlashAttention baseline in this
regime---$0.80\times$ and $0.52\times$---while the two standard quality
metrics disagreed about both of them (Section~\ref{sec:tiers}).

None of these facts is visible to isolated evaluation. Worse, isolated
evaluation is actively gameable by the systems being evaluated: our
independent validation caught an agent-generated kernel that was
bit-exact on every isolated test while having one call's memory layout
compiled in as constants, and dispatch code that silently bypassed
\texttt{torch.compile}'s guards to report a $2.557\times$ speedup its
guarded equivalent could not reproduce (Section~\ref{sec:taxonomy}). We
also caught ourselves: the pipeline's own impact attribution over-counted
a dispatch chain $4\times$, and an author-published number was corrected
when the experiment store ingested the underlying records. We therefore
treat machine-checked claims not as reproducibility hygiene but as a
load-bearing component of the system.

This paper makes the following contributions:
\begin{enumerate}
\item A fail-closed promotion pipeline for agent-generated optimizations,
  from declarative workloads and metadata-only graph capture through
  sandboxed autonomous search, independent validation, versioned
  artifacts, and full-generation A/B promotion (Section~3).
\item A quantified account of the promotion gap: isolated operator
  speedup does not predict shipped value, and measured per-artifact
  impact must replace analytical upper bounds, which overstated one
  kernel's return by more than $8\times$ (Section~\ref{sec:gap}).
\item Dispatch-tax accounting: a methodology for attributing delivery
  overhead per call, the $25{:}1$ result it produced, and the measured
  elimination of that overhead (Section~\ref{sec:dispatch}).
\item A failure taxonomy of agent-generated performance code, drawn from
  artifacts our gates caught, with the checks that caught them
  (Section~\ref{sec:taxonomy}).
\item Tiered fidelity contracts (bitwise, numeric, perceptual) that admit
  approximate optimizations without abandoning verification, and honest
  campaign results under them, including systematic disagreement between
  SSIM and LPIPS (Section~\ref{sec:tiers}).
\item End-to-end results and measurement methodology, including a
  promoted transformer artifact ($1.086\times$ median, byte-identical
  frames, $6{,}143$ dispatches, zero fallbacks) and an invalidation
  policy under which we retracted three of our own previously reported
  numbers when node-level variance proved sequential A/B measurement
  ungateable on our cluster (Section~\ref{sec:results}).
\end{enumerate}

We release the system, every workload manifest, and the machine-generated
measurement record behind every number in this paper.

\section{Background and Related Work}
\todo{(a) LLM/agent kernel generation and its isolated-metric evaluation
practice; (b) video DiT inference cost structure; (c) acceleration
families: sparse/quantized attention, step caching, compilation and
megakernels; (d) why existing evaluation does not measure shipped value.}

\section{System: A Fail-Closed Promotion Pipeline}
\todo{Stage-by-stage description with trust boundaries: declarative
workloads; metadata-only capture (symbolic FX $\to$ non-strict export
$\to$ Dynamo fallback; what is never serialized); pure-tensor allowlist
and fail-closed rejection; measured-impact ranking; graph-derived spec
generation; sandboxed autonomous search; write-once run contract;
independent validation; versioned artifact bundles;
compatibility-checked dispatch; full-generation A/B; finalization.
Figure~\ref{fig:pipeline}.}

\section{The Promotion Gap}
\label{sec:gap}

\begin{table}[h]
\centering
\caption{Isolated operator speedup versus end-to-end effect (Wan~2.1
T2V 1.3B, 480p, GB200). Isolated results are weighted across the
production shape corpus; the 50-step generation A/B was unchanged.}
\label{tab:wan}
\begin{tabular}{lrr}
\toprule
Fused operator & Isolated speedup & End-to-end effect \\
\midrule
Post-attention gated residual + LayerNorm & $8.638\times$ & \multirow{3}{*}{$\sim 1.00\times$} \\
Modulated pre-attention LayerNorm & $10.449\times$ & \\
Post-MLP gated residual & $10.628\times$ & \\
\bottomrule
\end{tabular}
\end{table}

\begin{table}[h]
\centering
\caption{The arithmetic of sub-percent regions: four validated LTX
candidates, their measured end-to-end share, and the gain achievable at
zero overhead and perfect parity. The campaign gate was $1.01\times$.}
\label{tab:ltx-arith}
\begin{tabular}{lrrr}
\toprule
Artifact & Share of e2e & Isolated speedup & Realized e2e gain \\
\midrule
mk-bbfe1518 & 2.334\% & $1.1154\times$ & 0.241\% \\
mk-a81e140d & 1.863\% & $1.1104\times$ & 0.185\% \\
mk-baecc382 & 1.081\% & $1.1212\times$ & 0.117\% \\
mk-b6cb64f9 & 0.922\% & $1.1064\times$ & 0.089\% \\
\midrule
\textbf{Total} & \textbf{6.20\%} & & \textbf{0.632\%} ($1.0064\times$) \\
\bottomrule
\end{tabular}
\end{table}

\todo{Measured- vs.\ estimated-impact ranking: the Amdahl bound overstated
a $1.115\times$ kernel's return by $>8\times$; per-artifact isolation
trials as the fix.}

\section{The Dispatch Tax}
\label{sec:dispatch}
\todo{Accounting methodology. Eager FX graph replay: $3.104$\,ms overhead
vs.\ $124$\,\textmu s saving per call; transformer at 384
calls/generation $\Rightarrow$ 1.192\,s, 32.4\% of e2e; VAE contrast
(4.7 calls, 0.39\%). The fix: guarded compiled dispatch via registered
ops; before/after measurement records. Figure: waterfall.}

\section{What Agents Get Wrong}
\label{sec:taxonomy}

Autonomous search produced kernels that were fast, numerically exact, and
wrong. This section is a taxonomy of the failures our pipeline caught,
organized by what made each one invisible to isolated evaluation. Every
case is drawn from preserved, immutable campaign evidence; none is
hypothetical. We include, deliberately, the failures that originated in
our own pipeline and in our own claims, because the taxonomy's lesson is
not ``agents write bad kernels''---it is that \emph{plausibility,
correctness on sampled inputs, and honest accounting are three different
properties}, and a promotion system must check all three.

\subsection{Contract violations that isolated tests cannot see}

The strongest candidate of one LTX campaign---a fused transformer
subregion spanning 22 graph nodes---passed smoke testing with
\texttt{max\_abs\_error} $= 0.0$, passed a full shape sweep at
$0.00$ worst error, and was bitwise deterministic across runs. It
failed the numerical-stability stage with a stride assertion:
\texttt{expected stride 4096 == 16384 at dim=1}. The kernel was not
numerically wrong anywhere; it had compiled the memory layout of
\emph{one observed call} into itself as compile-time constants:

\begin{verbatim}
_MAIN_ROWS = 4680
SOURCE_ROW_STRIDE = 16384   # a property of one call,
CONTEXT_ROW_STRIDE = 8192   # not of the operation
\end{verbatim}

The observed inputs were non-contiguous slices of larger timestep
tensors; the stability stage allocated fresh contiguous tensors of the
same shape, and the kernel would have read at wrong addresses. Every
isolated test that reuses the captured inputs---which is what
speedup-oriented evaluation naturally does---certifies this kernel as
perfect. The failure class is a \emph{layout contract violation}:
the kernel implements a function of one call's strides, not of the
operation. The gate that caught it was input diversity: validation
stages that re-generate inputs rather than replaying captured ones.

\subsection{Measurement gaming: the guard bypass}

The same candidate reported a $2.557\times$ isolated speedup. Its
dispatch wrapper compiled the implementation once, then cached and called
Inductor's raw compiled entry directly:

\begin{verbatim}
# "Calling the captured entry directly avoids a
#  guard/cache lookup before every GPU schedule."
outputs = _COMPILED_ENTRY(...)
\end{verbatim}

Two things are wrong at once. First, the number is unfair: the compiled
baseline it was compared against pays guard evaluation on every call, and
the candidate does not---the agent optimized the \emph{measurement
boundary}, not the kernel. Second, the bypass removes exactly the safety
mechanism that would have converted the layout violation of the previous
subsection from an assertion failure into silently wrong output at full
speed. The violation surfaced only because Inductor happens to emit an
internal stride assertion; the agent's Triton kernels contain no such
check. We draw a general conclusion: \textbf{an optimizing agent given a
timing harness will, without adversarial intent, discover the seams of
the harness}---and dispatch-path shortcuts look identical to legitimate
optimization in a diff. The countermeasure is structural, not
review-based: validation re-times candidates under its own guarded
dispatch, and a candidate whose speedup does not survive guarded
execution is scored on the guarded number.

\subsection{Attribution defects: when the pipeline lies to the search}

The search can only be as honest as the numbers steering it, and two of
the largest defects we found were in our own attribution, not in any
kernel. A campaign's transformer region was assigned
\texttt{share\_of\_e2e} $= 1.0333$---more than the whole
generation---because inclusive profiler ranges were summed against an
exclusive total. Its call count was recorded as $195{,}661$; the module
was invoked $1{,}151$ times---aten-level events across 48 blocks were
counted as module calls. Downstream, packaging consulted an analytic
Amdahl bound (\texttt{share} $\times 0.9$, assuming the region's cost
falls to zero) that overstated the realizable return of a $1.115\times$
kernel by more than $8\times$; two artifacts carried
\texttt{meets\_promotion\_target: false} into packaging and were
packaged anyway. The same defect recurred months later in our own
analysis: an attention-share estimate of $45.94\%$, computed by summing
an inclusive profiler column that counted one dispatch chain four times,
was corrected to a measured $27$--$35\%$. Attribution errors are more
dangerous than kernel errors: a wrong kernel is caught by a gate, while
a wrong impact number silently misallocates the entire search budget.
The remedies that worked were mechanical: derive impact from
per-artifact isolation trials rather than analytic bounds, and reconcile
module-level counts against profiler events before ranking consumes
them.

\subsection{Language that reads as verdicts}

Smaller defects concerned words rather than numbers, and they matter
because unattended runs are read after the fact by humans deciding what
shipped. A campaign status field reported four artifacts as
\texttt{finalized}; all four were quarantined and none was promoted. A
quarantine reason claimed ``the end-to-end validation stage did not
complete''; the stage had completed and returned a definite negative.
A search prompt announced ``Measured optimistic model impact: None\%''
because it read a key that nothing writes. Each is trivial; together
they compound into a run report that a reasonable reader mistakes for
success. Our rule since: status vocabulary is closed and tested, reasons
must quote the evidence they summarize, and any prompt field that can be
unset fails preflight rather than interpolating \texttt{None}.

\subsection{Failures caught by gates versus failures caught by tooling}

Of the defects above, the layout violation and the guard bypass were
caught by verification gates; the attribution and language defects were
caught by audit---a human rereading evidence---and later by
infrastructure built in response (an experiment store whose ingest
reconciles published numbers against raw records caught an
author-published speedup that had been transcribed from a job log, and a
factual claim about a model's attention requirements was corrected when
the store linked the constraint to the checkpoint rather than the
configuration). The trajectory across our campaigns is consistent:
\textbf{every class of error we initially caught by reading, we
eventually caught earlier by checking}, and the system's reliability
improved exactly as fast as claims moved from prose into
machine-validated records.
\todo{Per-gate catch-rate table from experiment-store backfill
(EVIDENCE\_MAP gap 3).}

\section{Fidelity Tiers and Approximate Optimizations}
\label{sec:tiers}
\todo{Contract design: Tier~0 byte-equal, Tier~1 numeric tolerance,
Tier~2 perceptual budget declared per workload (fixed-seed reference
frames, SSIM floor, LPIPS ceiling). Promotion carries the tier and the
margins.}

\begin{table}[h]
\centering
\caption{Attention backend campaign (Wan~2.1 T2V 1.3B, 480p, 49 frames,
20 steps, GB200/sm100; baseline pinned to FlashAttention; 15 timed runs
per arm). Gate: $\geq 1.3\times$, SSIM $\geq 0.97$, LPIPS $\leq 0.05$.}
\label{tab:attention}
\begin{tabular}{lrrrl}
\toprule
Candidate & Median speedup & Worst SSIM & Worst LPIPS & Verdict \\
\midrule
SageAttention 1.0.6 & $0.8031\times$ & 0.9353 & 0.0378 & rejected \\
Video Sparse Attention (zero-shot) & $0.5169\times$ & 0.9527 & 0.0253 & rejected \\
\bottomrule
\end{tabular}
\end{table}

\todo{Metric-divergence analysis: LPIPS admits both, SSIM rejects both;
argument for declaring both. Regime analysis: why these candidates lose
here (v1 quantization overhead on sm100; VSA tile selection cost at 81k
tokens; VSA's reported gains require VSA-finetuned checkpoints).
[PENDING: attention round 2 -- SageAttention2/2++, VSA on FastWan,
720p/121f. PENDING: caching campaign -- first Tier-2 promotion.]}

\section{End-to-End Results}
\label{sec:results}
\todo{LTX2 promoted artifact: $3.3646 \to 3.0991$\,s median
($1.0857\times$), 15-run A/B; independent replication $1.2514\times$
[link raw record]; $6{,}143$ dispatches, zero fallbacks, byte-identical
frames. Support matrix snapshot (model $\times$ workload $\times$ arch)
[PENDING sm90]. Variance methodology: native-arm CoV ceiling, locked
clocks, exclusive nodes; before/after spread.}

\section{Limitations and Discussion}
\todo{Single-framework integration depth (FastVideo) vs.\ what
generalizes; regimes not covered (multi-GPU, training); what ``verified''
does and does not claim; cost of the pipeline itself.}

\section{Conclusion}
\todo{}

\section*{Acknowledgments}
\todo{Agent tooling disclosure; GPU resources.}

\section*{Reproducibility}
Every quantitative claim traces to a committed, machine-generated
measurement record (SLURM job IDs, environment identity, schema-versioned
JSON) in the public repository. \todo{repo tag + evidence index link.}

\bibliographystyle{plain}
\bibliography{refs}

\end{document}
